\specialchapter{历史注记}

（注意 --- 罗马数字指本注记末尾的参考文献。）

尽管拓扑学与测度论之间联系的研究可以追溯到现代实变量函数论的开端，但在 Hausdorff 拓扑空间中的积分以一般的方式得到彻底研究，还只是最近的事。在叙述先于本综合工作的那些工作的历史之前，我们回顾关于拓扑与测度之间关系的观念演化的几个阶段。

对 Lebesgue 而言，问题只在于积分一个或几个实变量的函数。1913 年，Radon 定义了 \(\mathbf{R}^n\) 上的一般测度以及相应的积分；这一理论在 Ch. de la Vallée Poussin 的书 (I) 中有详细的阐述，并且自始至终依赖于 Euclid 空间的拓扑性质。稍后，在 1915 年，Fréchet 在 (II, \(a\) ) 中定义了赋有一个 σ-集环的集合上的“抽象”测度以及关于这些测度的积分；他注意到，用这种方式可以不借助拓扑方法而建立 Lebesgue 理论的主要结果。他用取自 (II, \(b\) ) 引言（第一部分）的下面这段话来为自己的工作辩护：« Que par exemple dans l'espace à une infinité de coordonnées où diverses applications de l'Analyse avaient conduit à diverses définitions non équivalentes d'une suite convergente, on remplace une de ces définitions par une autre, rien ne sera changé dans les propriétés des familles et fonctions additives d'ensembles dans ces espaces».(*) Fréchet 的研究由 Carathéodory 完成，集函数扩张为测度的一个重要定理应归功于后者。Saks 的书 (III) 的开头对这一观点作了扼要的阐述。

Haar 测度在局部紧群上的发现（cf.~第七章与第八章的历史注记）以及它随即获得的众多应用，随后 Weil 与 Gelfand 在调和分析方面的工作，大约在 1940 年前后导致了这一观点的深刻改变：在这类问题中，把测度视为连续函数空间上的线性形式最为方便。这一方法使人不得不限于紧空间或局部紧空间，但对几乎所有的应用而言这并无不便：更好的是，J. Tate 与 A. Weil 在 p-adic 群与 adèle 群上引入调和分析，使解析数论获得了惊人的更新。

正是从一个完全不同的方向，产生了通过考虑非局部紧拓扑空间上的测度来扩大这一观点的需要：概率论逐渐引导到对这类空间的研究，并提供了许多非平凡的例子。这些发展对测度论的影响之所以姗姗来迟，其原因或许在于概率论的相对孤立，直到不久以前它还处于传统数学学科的边缘。

